%----------------------------------------------------------------------------------------
%  Geopolitics for Business — Group Project Presentation
%  Data Group 2
%  Template: layout.tex (Darmstadt / Bocconi blue)
%----------------------------------------------------------------------------------------

\documentclass[10pt]{beamer}

\mode<presentation> {
  \usetheme{Darmstadt}

  % Remove top navigation strip
  \setbeamertemplate{headline}{}

  % Bocconi blue palette (replaces seahorse)
  \definecolor{bocblue}{HTML}{002B5B}
  \setbeamercolor{structure}{fg=bocblue}
  \setbeamercolor{palette primary}{bg=bocblue,fg=white}
  \setbeamercolor{palette secondary}{bg=bocblue!80,fg=white}
  \setbeamercolor{palette tertiary}{bg=bocblue!60,fg=white}
  \setbeamercolor{palette quaternary}{bg=bocblue,fg=white}
  \setbeamercolor{titlelike}{parent=palette primary}
  \setbeamercolor{frametitle}{bg=bocblue,fg=white}
  \setbeamercolor{block title}{bg=bocblue,fg=white}
  \setbeamercolor{block body}{bg=bocblue!10}
  \setbeamercolor{section in head/foot}{bg=bocblue,fg=white}
  \setbeamercolor{subsection in head/foot}{bg=bocblue!70,fg=white}
  \setbeamertemplate{blocks}[rounded][shadow=false]
}

\usepackage{graphicx}
\usepackage{booktabs}
\usepackage{colortbl,xcolor}
\usepackage{amsmath,amssymb}
\usepackage{hyperref}

% Graceful placeholder when figure files are missing
\newcommand{\figureOrBox}[2][0.85\linewidth]{%
  \IfFileExists{#2}{%
    \includegraphics[width=#1]{#2}%
  }{%
    \fbox{\parbox{#1}{\centering\vspace{1.2em}\texttt{\tiny #2}\\[0.3em]%
      \textit{\tiny Place figure here}\vspace{1.2em}}}%
  }%
}

%----------------------------------------------------------------------------------------
%  TITLE PAGE
%----------------------------------------------------------------------------------------

\title[Geopolitics for Business]{Energy Shock and Firm-Level Decarbonisation}
\subtitle{Did the Russia--Ukraine War Reduce GHG Intensity in European Industry?}

\author{Data Group 2}
\institute[Bocconi University]{
  Marius Borsu Lefevre \quad Giovanni Clemente \quad Greta Guadalupi \\
  Petr Hrebenar \quad Stéphane Miquel Clément \quad Pietro Montano \\
  Souhail Habib Nehme \quad Joan Olaf Subirachs Garcia \\[0.8em]
  \small Bocconi University --- Geopolitics for Business
}
\date{May 2026}

\begin{document}

\begin{frame}
  \titlepage
\end{frame}

%----------------------------------------------------------------------------------------
%  SLIDE 2: BIG PICTURE
%----------------------------------------------------------------------------------------
\section{Motivation}

\begin{frame}{Big Picture}
  \begin{block}{The context}
    Russia's 2022 invasion of Ukraine triggered the largest European energy supply disruption in decades,
    sending gas prices soaring and forcing firms to rethink their energy sourcing overnight.
  \end{block}

  \vspace{0.6em}

  \begin{itemize}
    \item Europe had built production models around cheap Russian pipeline gas
    \item The shock arrived as climate commitments tightened: EU Fit for 55, CSRD
    \item Energy-intensive industries --- chemicals, metals, cement, paper --- were hit hardest
    \item Yet the literature on geopolitical shocks and environmental performance is thin
  \end{itemize}
\end{frame}

%----------------------------------------------------------------------------------------
%  SLIDE 3: WHY WE CARE
%----------------------------------------------------------------------------------------

\begin{frame}{Why We Care}
  \begin{block}{The gap}
    We know a lot about how geopolitical shocks affect investment and profitability.
    We know far less about their effect on \textbf{firm-level environmental performance}.
  \end{block}

  \vspace{0.6em}

  \begin{itemize}
    \item Quantifying the firm-level GHG response is essential to evaluate whether
          energy-security shocks and climate policy \emph{reinforce} or \emph{undermine} each other
    \item Country-level energy infrastructure may become a \emph{bottleneck} for decarbonisation ---
          even when price incentives are strong
    \item Policy implication: energy-mix flexibility at the country level enables (or blocks)
          firm-level environmental adjustment
  \end{itemize}
\end{frame}

%----------------------------------------------------------------------------------------
%  SLIDE 4: RESEARCH QUESTION
%----------------------------------------------------------------------------------------
\section{Research Question}

\begin{frame}{Research Question}
  \begin{block}{}
    \normalsize
    Did the Russia--Ukraine energy shock trigger a \textbf{reduction in firm-level GHG intensity}
    among European energy-intensive firms --- and was this effect \textbf{different} for firms
    in countries more or less dependent on Russian gas?
  \end{block}

  \vspace{0.8em}

  \begin{itemize}
    \item \textbf{Outcome:} GHG Direct Intensity --- Scope 1 emissions (tCO$_2$e) per \$1M revenue
    \item Measured by \textbf{TruCost (S\&P Global)}: captures \emph{operational} emissions efficiency,
          normalised for firm size, with lowest measurement noise across datasets (NGFS 2024)
  \end{itemize}
\end{frame}

%----------------------------------------------------------------------------------------
%  SLIDE 5: DATA
%----------------------------------------------------------------------------------------
\section{Data}

\begin{frame}{Data}
  \textbf{Panel dataset}
  \begin{itemize}
    \item $\approx$973{,}000 firm-year observations
    \item 28 European countries (EU-27 + NO, CH, UK)
    \item Years: 2018--2024 (pre-shock: 2018--2021; post-shock: 2022--2024)
    \item Publicly listed firms in manufacturing and resource sectors
  \end{itemize}

  \vspace{0.6em}

  \textbf{Sources}
  \begin{itemize}
    \item \textbf{TruCost (S\&P Global)} --- GHG Direct Intensity (Scope 1)
    \item \textbf{World Bank} --- GDP per capita, CPI inflation
    \item \textbf{Eurostat} --- Russian gas import shares (\texttt{nrg\_ti\_gas}, \texttt{nrg\_bal\_c})
    \item \textbf{EMAPS / Ember} --- Country-year electricity prices
  \end{itemize}
\end{frame}

%----------------------------------------------------------------------------------------
%  SLIDE 6: METHODOLOGY 1 — DiD IDENTIFICATION STRATEGY
%----------------------------------------------------------------------------------------
\section{Methodology}

\begin{frame}{Identification Strategy: Difference-in-Differences}
  \begin{itemize}
    \item \textbf{Core idea:} compare the change in GHG intensity for \emph{treated} firms
          relative to \emph{control} firms, before vs.\ after February 2022
    \item \textbf{Why the shock is exogenous:} the RU-UA war and energy disruptions were
          unanticipated by firms --- no pre-announcement or anticipation bias
    \item \textbf{Parallel trends:} treated and control firms show broadly parallel
          GHG intensity trends in 2018--2021; validated via event-study pre-period coefficients
  \end{itemize}

  \vspace{0.4em}

  \begin{figure}
    \centering
    \figureOrBox[0.90\linewidth]{final_output/parallel_trends_plot.png}
    \caption{\footnotesize Mean GHG Direct Intensity, treated vs.\ control. Dotted line: Feb 2022 shock.}
  \end{figure}
\end{frame}

%----------------------------------------------------------------------------------------
%  SLIDE 7: METHODOLOGY 2 — TREATMENT 1: GAS EXPOSURE
%----------------------------------------------------------------------------------------

\begin{frame}{Treatment 1: Country Gas Exposure}
  \begin{columns}[T]
    \begin{column}{0.42\textwidth}
      \small
      \[
        \text{Exp}_{c,t} = \frac{\text{RU gas imports}_{c,t}}{\text{Gross Inland Consumption}_{c,t}}
      \]
      \vspace{0.2em}
      \begin{itemize}\small
        \item Pre-shock average 2018--2021
        \item \textbf{Treated:} top quartile ($\geq$ p75 $=$ 14.1\%)
        \item 8 countries: HU, LV, SK, IT, CZ, DE, LT, BG
        \item Source: Eurostat \texttt{nrg\_ti\_gas}, \texttt{nrg\_bal\_c}
      \end{itemize}
    \end{column}
    \begin{column}{0.55\textwidth}
      \figureOrBox[\linewidth]{final_output/russian_exposure_barplot.png}
      \vspace{-0.3em}
      \begin{center}
        \tiny Avg.\ Russian gas exposure by country, 2018--2021. Red = treated ($\geq$ p75).
      \end{center}
    \end{column}
  \end{columns}
\end{frame}

%----------------------------------------------------------------------------------------
%  SLIDE 8: METHODOLOGY 3 — MODEL 1 SPECIFICATION
%----------------------------------------------------------------------------------------

\begin{frame}{Model 1: Gas Exposure Specification}
  \begin{equation}
    Y_{i,t} = \alpha_i + \lambda_t + \beta_1 \cdot \text{Post}_t \times \text{GasExp}_i + \gamma' X_{c,t} + \varepsilon_{i,t}
  \end{equation}

  \vspace{0.3em}

  \begin{itemize}
    \item $Y_{i,t}$: GHG Direct Intensity (tCO$_2$e / \$M revenue)
    \item $\alpha_i$: firm fixed effects \quad $\lambda_t$: year fixed effects
    \item $\text{GasExp}_i = 1$ if firm's country is in the top-p75 exposure group
    \item $\text{Post}_t = \mathbf{1}(t \geq 2022)$
    \item $X_{c,t}$: GDP per capita + CPI inflation (country-year controls)
    \item SE clustered at the \textbf{country level}
  \end{itemize}

  \vspace{0.4em}

  \begin{block}{Coefficient of interest}
    \small $\hat\beta_1$: avg.\ change in GHG intensity for gas-exposed countries vs.\ non-exposed, post-2022 vs.\ pre-2022
  \end{block}
\end{frame}

%----------------------------------------------------------------------------------------
%  SLIDE 9: METHODOLOGY 4 — TREATMENT 2: EII SECTORS
%----------------------------------------------------------------------------------------

\begin{frame}{Treatment 2: Energy-Intensive Sectors (EII)}
  \begin{columns}[T]
    \begin{column}{0.48\textwidth}
      \textbf{Treated (EII = 1)}\\[0.3em]
      OECD/EC official EII classification:
      \begin{itemize}\small
        \item C16 Wood products
        \item C17 Paper
        \item C19 Coke \& refined petroleum
        \item C20 Chemicals
        \item C22 Rubber \& plastics
        \item C23 Non-metallic minerals (cement, glass)
        \item C24 Basic metals
      \end{itemize}
    \end{column}
    \begin{column}{0.48\textwidth}
      \textbf{Control (EII = 0)}\\[0.3em]
      Comparable manufacturing, lower energy input share:
      \begin{itemize}\small
        \item C21 Pharmaceuticals
        \item C25 Fabricated metals
        \item C28 Machinery
        \item C29 Motor vehicles
        \item C31 Furniture
        \item C18 Printing
      \end{itemize}
      \vspace{0.5em}
      \small 66\% of observations in EII sectors (647K firm-years)
    \end{column}
  \end{columns}
\end{frame}

%----------------------------------------------------------------------------------------
%  SLIDE 10: METHODOLOGY 5 — MODEL 2 SPECIFICATION
%----------------------------------------------------------------------------------------

\begin{frame}{Model 2: Sector Energy Intensity Specification}
  \begin{equation}
    Y_{i,t} = \alpha_i + \lambda_t + \beta_2 \cdot \text{Post}_t \times \text{EII}_i + \gamma' X_{c,t} + \varepsilon_{i,t}
  \end{equation}

  \vspace{0.3em}

  \begin{itemize}
    \item $\text{EII}_i = 1$ for energy-intensive sectors (C16, C17, C19, C20, C22, C23, C24)
    \item Same firm/year FE and country-clustered SE as Model 1
  \end{itemize}

  \vspace{0.4em}

  \begin{block}{Coefficient of interest}
    \small $\hat\beta_2$: avg.\ change in GHG intensity for EII firms vs.\ non-EII, post-2022 vs.\ pre-2022
  \end{block}
\end{frame}

%----------------------------------------------------------------------------------------
%  SLIDE 11: METHODOLOGY 6 — FINAL MODEL (MODEL 4)
%----------------------------------------------------------------------------------------

\begin{frame}{Final Model: Triple Interaction (Model 3)}
  \small
  \begin{equation}
    Y_{i,t} = \alpha_i + \lambda_t
      + \beta_1 \, \text{Post}_t \times \text{GasExp}_i
      + \beta_2 \, \text{Post}_t \times \text{EII}_i
      + \beta_3 \, \text{Post}_t \times \text{GasExp}_i \times \text{EII}_i
      + \gamma' X_{c,t} + \varepsilon_{i,t}
  \end{equation}

  \normalsize\vspace{0.3em}

  \begin{itemize}
    \item $\hat\beta_1$: standalone country gas-exposure effect
    \item $\hat\beta_2$: standalone EII sector effect
    \item $\hat\beta_3$: \textbf{additional} effect for firms that are \emph{both}
          in gas-exposed countries \emph{and} in EII sectors
  \end{itemize}

  \vspace{0.3em}

  \begin{block}{Why triple interaction?}
    Allows us to decompose country- and sector-level channels separately,
    and test whether the two dimensions of treatment reinforce or offset each other.
  \end{block}
\end{frame}

%----------------------------------------------------------------------------------------
%  SLIDE 12: RESULTS — REGRESSION TABLE
%----------------------------------------------------------------------------------------
\section{Results}

\begin{frame}{Results: All Models}
  \begin{center}
    \small
    \begin{tabular}{lccc}
      \toprule
       & (1) & (2) & (3) \\
       & Gas Exposed & Energy Intensive & Triple \\
      \midrule
      Post $\times$ GasExp
        & $5.209^{***}$ & $-38.461^{***}$ & $-2.611$ \\
        & $(1.089)$ & $(2.209)$ & $(2.498)$ \\[3pt]
      Post $\times$ EII
        & & & $-43.480^{***}$ \\
        & & & $(1.643)$ \\[3pt]
      Post $\times$ Both
        & & & $9.863^{***}$ \\
        & & & $(2.543)$ \\[3pt]
      \midrule
      Observations & 972{,}979 & 972{,}979 & 972{,}979 \\
      Within $R^2$  & 0.011 & 0.048 & 0.049 \\
      Firm / Year FE & Yes & Yes & Yes \\
      SE clustered & Country & Country & Country \\
      \bottomrule
    \end{tabular}
    \par\vspace{0.4em}
    {\tiny Dep.\ var.: GHG Direct Intensity (tCO$_2$e / \$1M revenue).
      SE in parentheses. $^{*}p<0.10$, $^{**}p<0.05$, $^{***}p<0.01$.}
  \end{center}
\end{frame}

%----------------------------------------------------------------------------------------
%  SLIDE 13: KEY FINDING — MODEL 2
%----------------------------------------------------------------------------------------

\begin{frame}{Finding: EII Firms Reduced GHG Intensity}
  \begin{block}{Model 2}
    \small EII firms reduced GHG Direct Intensity by $\approx \mathbf{38.5}$\,tCO$_2$e/\$M post-2022 vs.\ non-EII firms\\
    \hfill ($p < 0.01$; $\approx$32\% of sample mean of 121\,tCO$_2$e/\$M)
  \end{block}

  \vspace{0.3em}

  \begin{itemize}
    \item Higher energy prices forced diversification away from gas $\to$ Scope 1 falls
    \item Effect is sustained across 2022--2024, not a one-off adjustment
  \end{itemize}

  \vspace{0.2em}

  \begin{figure}
    \centering
    \figureOrBox[0.85\linewidth]{final_output/event_study_plot.png}
    \caption{\footnotesize Event study: $\hat\beta_t$ for EII vs.\ non-EII, normalised to 2021.}
  \end{figure}
\end{frame}

%----------------------------------------------------------------------------------------
%  SLIDE 14: KEY FINDING — MODEL 4
%----------------------------------------------------------------------------------------

\begin{frame}{Finding: Gas Dependency Attenuates Decarbonisation}
  \begin{block}{Model 3}
    \small $\hat\beta_2 = -43.5^{***}$, $\hat\beta_3 = +9.8^{***}$ $\Rightarrow$ net effect (doubly treated): $\mathbf{-33.7}$\,tCO$_2$e/\$M
  \end{block}

  \vspace{0.3em}

  \begin{itemize}
    \item Gas-dependent countries faced \textbf{constrained transition paths} (LNG bottlenecks, long renewable lead times)
    \item Standalone gas-exposure effect: $+5.2$ in Model 1 $\to$ indistinguishable from zero ($-2.6$, $p=0.30$) in Model 3
  \end{itemize}

  \vspace{0.2em}

  \begin{figure}
    \centering
    \figureOrBox[0.78\linewidth]{final_output/event_study_plot_triple.png}
    \caption{\footnotesize Dynamic DiD: total effect for doubly-treated firms ($\hat\beta_{2,t}+\hat\beta_{3,t}$). Ref: 2021.}
  \end{figure}
\end{frame}

%----------------------------------------------------------------------------------------
%  SLIDE 15: DISCUSSION
%----------------------------------------------------------------------------------------
\section{Discussion}

\begin{frame}{Discussion: Mechanism}
  \begin{enumerate}
    \item Energy price shock $\to$ strong incentive to cut gas consumption
    \item Firms diversify: energy efficiency, lower-carbon electricity, supply-chain changes
    \item Gas is itself a GHG fuel $\to$ diversification \emph{mechanically} reduces Scope 1
  \end{enumerate}

  \vspace{0.5em}

  \begin{center}
    \small
    \begin{tabular}{lll}
      \toprule
      Country & Energy endowment & Adjustment capacity \\
      \midrule
      France & Nuclear flexibility & Higher \\
      Germany & Coal fallback + renewables & Moderate \\
      Italy & LNG ramp-up constrained & Lower \\
      HU / SK & High pipeline dependency & Lowest \\
      \bottomrule
    \end{tabular}
  \end{center}

  \vspace{0.3em}

  Firm-level decarbonisation and \textbf{country-level energy infrastructure} are deeply complementary.
\end{frame}

%----------------------------------------------------------------------------------------
%  SLIDE 16: LIMITATIONS
%----------------------------------------------------------------------------------------
\section{Limitations}

\begin{frame}{Limitations}
  \begin{enumerate}
    \item \textbf{HQ $\neq$ operations for MNEs}\\
      {\small A German-registered firm producing in Poland faces the Polish energy mix.
      Facility-level data would be needed to correct this.}

    \vspace{0.4em}

    \item \textbf{Carbon leakage not identified}\\
      {\small Firms may have outsourced emission-intensive steps abroad.
      Scope 3 or trade-adjusted accounting required.}

    \vspace{0.4em}

    \item \textbf{Pre-trend caveat (EII treatment)}\\
      {\small 2018 pre-period coefficient elevated ($\approx +16$); results capture \emph{acceleration}
      of a divergence, not its initiation. Post-shock drop ($-25$ to $-35$) is an order of magnitude larger.}

    \vspace{0.4em}

    \item \textbf{Short post-shock window}\\
      {\small Panel ends in 2024 --- only two full post-shock years. Full trajectory may materialise
      as renewable capacity investments complete.}
  \end{enumerate}
\end{frame}

%----------------------------------------------------------------------------------------
%  SLIDE 17: CONCLUSION + THANK YOU
%----------------------------------------------------------------------------------------
\section{Conclusion}

\begin{frame}{Conclusion}
  \begin{block}{Main finding}
    The RU-UA energy shock \textbf{significantly reduced GHG intensity} among European EII firms:
    $-38$--$43$\,tCO$_2$e/\$M (32--36\% of sample mean).
    Effect is \textbf{smaller} in gas-dependent countries ($\approx -10$\,tCO$_2$e/\$M attenuation).
  \end{block}

  \vspace{0.5em}

  \textbf{Policy implications}

  \begin{columns}[T]
    \begin{column}{0.48\textwidth}
      \textbf{Country level}\\[0.2em]
      {\small Energy diversification (interconnectors, LNG, renewables) is both an energy-security
      \emph{and} decarbonisation imperative.}
    \end{column}
    \begin{column}{0.48\textwidth}
      \textbf{Firm level}\\[0.2em]
      {\small Prioritise energy-mix flexibility and long-term power-purchase agreements
      to decouple GHG performance from host-country endowments.}
    \end{column}
  \end{columns}

\end{frame}

\end{document}
